
\documentclass{article}
\usepackage{colortbl}
\usepackage{makecell}
\usepackage{multirow}
\usepackage{supertabular}

\begin{document}

\newcounter{utterance}

\centering \large Interaction Transcript for game `referencegame', experiment `line\_grids\_columns', episode 5 with reflexion{-}Qwen3.5{-}9B.
\vspace{24pt}

{ \footnotesize  \setcounter{utterance}{1}
\setlength{\tabcolsep}{0pt}
\begin{supertabular}{c@{$\;$}|p{.15\linewidth}@{}p{.15\linewidth}p{.15\linewidth}p{.15\linewidth}p{.15\linewidth}p{.15\linewidth}}
    \# & $\;$A & \multicolumn{4}{c}{Game Master} & $\;\:$B\\
    \hline

    \theutterance \stepcounter{utterance}  
    & & \multicolumn{4}{p{0.6\linewidth}}{
        \cellcolor[rgb]{0.9,0.9,0.9}{
            \makecell[{{p{\linewidth}}}]{
                \texttt{\tiny{[P1$\langle$GM]}}
                \texttt{You have attempted this same task before and did not succeed. These are the plans you wrote after those attempts.} \\
\texttt{Trial \#1: New plan: I incorrectly identified the columns as the first and fourth, but the target grid clearly shows the first and last columns (columns 1 and 5) are filled with 'X', while the middle three columns are empty. My strategy failed because I miscounted the column indices. Next time, I will carefully verify the column positions by counting from left to right (1 to 5) before describing the pattern. I will explicitly state that the first and fifth columns are filled to ensure the description uniquely distinguishes the target from the distractors, which have different column fill patterns.} \\
\\ 
\\ 
\texttt{You are given three grids, where each of them is 5 by 5 in size.} \\
\texttt{Grids have empty cells marked with "▢" and filled cells marked with "X".} \\
\texttt{Your task is to generate a referring expression that best describes the target grid while distinguishing it from the two other distractor grids.} \\
\texttt{The first grid is the target grid, and the following two grids are the distractors.} \\
\\ 
\texttt{Target grid:} \\
\\ 
\texttt{X □ □ X X} \\
\texttt{X □ □ X X} \\
\texttt{X □ □ X X} \\
\texttt{X □ □ X X} \\
\texttt{X □ □ X X} \\
\\ 
\texttt{Distractor grid 1:} \\
\\ 
\texttt{□ □ X X X} \\
\texttt{□ □ X X X} \\
\texttt{□ □ X X X} \\
\texttt{□ □ X X X} \\
\texttt{□ □ X X X} \\
\\ 
\texttt{Distractor grid 2:} \\
\\ 
\texttt{X □ X X □} \\
\texttt{X □ X X □} \\
\texttt{X □ X X □} \\
\texttt{X □ X X □} \\
\texttt{X □ X X □} \\
\\ 
\texttt{Instruction: Describe the target grid.} \\
\texttt{Generate the referring expression starting with the tag "Expression: " for the given target grid. Omit any other text.} \\
            }
        }
    }
    & \\ \\

    \theutterance \stepcounter{utterance}  
    & \multicolumn{4}{p{0.6\linewidth}}{
        \cellcolor[rgb]{0.9,0.9,0.9}{
            \makecell[{{p{\linewidth}}}]{
                \texttt{\tiny{[P1$\rangle$GM]}}
                \texttt{Expression: The grid where the first and last columns are filled with 'X' and the middle three columns are empty.} \\
            }
        }
    }
    & & \\ \\

    \theutterance \stepcounter{utterance}  
    & & & \multicolumn{2}{p{0.3\linewidth}}{
        \cellcolor[rgb]{0.9,0.9,0.9}{
            \makecell[{{p{\linewidth}}}]{
                \texttt{\tiny{[GM$|$GM]}}
                \texttt{Expression: The grid where the first and last columns are filled with 'X' and the middle three columns are empty.} \\
            }
        }
    }
    & & \\ \\

    \theutterance \stepcounter{utterance}  
    & & \multicolumn{4}{p{0.6\linewidth}}{
        \cellcolor[rgb]{0.9,0.9,0.9}{
            \makecell[{{p{\linewidth}}}]{
                \texttt{\tiny{[GM$\rangle$P2]}}
                \texttt{You have attempted this same task before and did not succeed. These are the plans you wrote after those attempts.} \\
\texttt{Trial \#1: New plan: I incorrectly identified the "First" grid because I focused on the vertical columns of Xs without verifying the horizontal rows. The expression requires the *first and fourth columns* to be filled, meaning every cell in column 1 and column 4 must be "X", while columns 2, 3, and 5 must be empty. In the "First" grid, the first column is filled, but the fourth column contains "X"s in the first two rows and empty cells below, which contradicts the requirement that the fourth column be fully filled. The "Second" grid has the first two columns filled, not the first and fourth. The "Third" grid has the first column filled, but the fourth column is empty. Upon re{-}evaluating the grids against the strict column constraints, I realize none of the provided grids perfectly match the description of having *both* the first and fourth columns entirely filled with Xs while the others are empty. However, if I must choose the closest match based on the primary feature (first column filled), I should have scrutinized the fourth column more closely. Actually, looking closer at the "First" grid: Col 1 is all X, Col 4 is X,X,X,X,▢ (not all X). The "Second" grid: Col 1 is all ▢. The "Third" grid: Col 1 is all X, Col 4 is all X. Wait, let me re{-}read the Third grid: "X □ □ X X". Here, Col 1 is X, Col 2 is □, Col 3 is □, Col 4 is X, Col 5 is X. This matches "first and fourth columns are filled" (Col 1 and Col 4 are X) but the expression says "remaining three columns are empty". In the Third grid, the remaining columns are 2, 3, and 5. Col 5 is filled with X. So the Third grid fails the "remaining three columns are empty" condition. The First grid has Col 5 empty, but Col 4 is not fully filled. The Second grid has Col 1 empty. It seems there is no perfect match, or I misread the grid layout. Let's re{-}examine the "First" grid: "X □ X X □". Col 1: X. Col 2: □. Col 3: X. Col 4: X. Col 5: □. This grid has Xs in cols 1, 3, 4. The expression wants only cols 1 and 4 filled. The "Third" grid: "X □ □ X X". Cols 1, 4, 5 are filled. The expression wants only cols 1 and 4. Neither grid perfectly matches. However, the "First" grid has the most Xs in the correct positions (1 and 4} \\
\\ 
\\ 
\texttt{You are given three grids, where each of them is 5 by 5 in size.} \\
\texttt{Grids have empty cells marked with "▢" and filled cells marked with "X".} \\
\texttt{You are also given a referring expression that describes one of the given grids.} \\
\texttt{Your task is to select the grid that matches the given referring expression.} \\
\texttt{Generate only the number (in text) of the grid that matches the given expression by selecting first, second, or third.} \\
\\ 
\texttt{First:} \\
\\ 
\texttt{X □ X X □} \\
\texttt{X □ X X □} \\
\texttt{X □ X X □} \\
\texttt{X □ X X □} \\
\texttt{X □ X X □} \\
\\ 
\texttt{Second:} \\
\\ 
\texttt{□ □ X X X} \\
\texttt{□ □ X X X} \\
\texttt{□ □ X X X} \\
\texttt{□ □ X X X} \\
\texttt{□ □ X X X} \\
\\ 
\texttt{Third:} \\
\\ 
\texttt{X □ □ X X} \\
\texttt{X □ □ X X} \\
\texttt{X □ □ X X} \\
\texttt{X □ □ X X} \\
\texttt{X □ □ X X} \\
\\ 
\texttt{Expression: The grid where the first and last columns are filled with 'X' and the middle three columns are empty.} \\
\texttt{Question: Which grid does the expression refer to?} \\
\texttt{Start with the tag "Answer: ", followed by your selection. Omit any other text.} \\
            }
        }
    }
    & \\ \\

    \theutterance \stepcounter{utterance}  
    & & & \multicolumn{4}{p{0.6\linewidth}}{
        \cellcolor[rgb]{0.9,0.9,0.9}{
            \makecell[{{p{\linewidth}}}]{
                \texttt{\tiny{[GM$\langle$P2]}}
                \texttt{Answer: first} \\
            }
        }
    }
     \\ \\

    \theutterance \stepcounter{utterance}  
    & & & \multicolumn{2}{p{0.3\linewidth}}{
        \cellcolor[rgb]{0.9,0.9,0.9}{
            \makecell[{{p{\linewidth}}}]{
                \texttt{\tiny{[GM$|$GM]}}
                \texttt{first} \\
            }
        }
    }
    & & \\ \\

    \theutterance \stepcounter{utterance}  
    & \multicolumn{4}{p{0.6\linewidth}}{
        \cellcolor[rgb]{0.95,0.95,0.95}{
            \makecell[{{p{\linewidth}}}]{
                \texttt{\tiny{[P2$\rightarrow$P1]}}
                \texttt{Post{-}episode Reflexion (not sent to GM):} \\
\texttt{New plan: My previous description was ambiguous because it failed to explicitly mention the "middle three columns" being empty, which is the critical feature distinguishing the target from Distractor 1 (where the rightmost columns are filled) and Distractor 2 (where the pattern alternates). To ensure uniqueness, I must describe the target as having the first and last columns filled with 'X' while explicitly stating that the three central columns (columns 2, 3, and 4) are entirely empty. This precise phrasing will clearly separate the target grid from both distractors.} \\
\\ 
\texttt{Expression: The grid where the first and last columns are filled with 'X' and the middle three columns are empty.} \\
            }
        }
    }
    & & \\ \\

    \theutterance \stepcounter{utterance}  
    & & & \multicolumn{4}{p{0.6\linewidth}}{
        \cellcolor[rgb]{0.95,0.95,0.95}{
            \makecell[{{p{\linewidth}}}]{
                \texttt{\tiny{[P1$\rightarrow$P2]}}
                \texttt{Post{-}episode Reflexion (not sent to GM):} \\
\texttt{New plan: I failed because I misidentified the column indices of the "First" grid, incorrectly assuming the pattern matched the description when it actually had filled cells in columns 1, 3, and 4. The expression strictly requires only the first and last columns (1 and 5) to be filled, with columns 2, 3, and 4 empty. I will re{-}evaluate each grid by explicitly checking the content of every column index: verify that column 1 and column 5 contain only "X", and columns 2, 3, and 4 contain only "▢". Upon this strict check, I will select the grid that satisfies all five column constraints simultaneously.} \\
            }
        }
    }
     \\ \\

\end{supertabular}
}

\end{document}
